% =============================================================================
\section{Tarefa 3 --- Inclinómetro com o acelerómetro}
\label{sec:incl}
% =============================================================================

Se o sensor não acelera, o acelerómetro calibrado mede só a gravidade, e a
sua direção no referencial do sensor dá a inclinação. Com
$\vect a = (a_x;\,a_y;\,a_z)$ a medição corrigida,
\begin{equation}
\phi = \operatorname{atan2}(a_y,\ a_z), \qquad
\theta = \operatorname{atan2}\!\bigl(-a_x,\ \sqrt{a_y^2 + a_z^2}\bigr),
\label{eq:incl}
\end{equation}
onde $\phi$ é o rolamento e $\theta$ a arfagem. A rotação em torno da
vertical não altera a gravidade vista pelo sensor, por isso não se obtém.

O método tem dois limites. É ruidoso, porque cada estimativa usa uma só
amostra. E assume aceleração nula: qualquer aceleração do sensor é lida como
inclinação. O módulo da medição indica quando isso acontece, porque se afasta
de $1\,g$.

\todo{resultados no ensaio D (erro médio e máximo contra o ângulo do robô) e
no D2 (erro durante a perturbação); remeter para a Figura~\ref{fig:fusao}.}
